\documentclass[10pt,a4paper]{article}
\usepackage{amsmath}
\usepackage{fontspec}
\usepackage{babel}
\usepackage[margin=2.5cm]{geometry}
\usepackage{enumitem}
\usepackage{xcolor}
\usepackage{hyperref}
\usepackage{mdframed}

%% ─── Design system tokens (@docs/design/design.md) ──────────
\definecolor{primary}{HTML}{cc785c}
\definecolor{coral}{HTML}{cc785c}
\definecolor{primary-active}{HTML}{a9583e}
\definecolor{primary-disabled}{HTML}{e6dfd8}
\definecolor{ink}{HTML}{141413}
\definecolor{body-strong}{HTML}{252523}
\definecolor{body}{HTML}{3d3d3a}
\definecolor{muted}{HTML}{6c6a64}
\definecolor{muted-soft}{HTML}{8e8b82}
\definecolor{canvas}{HTML}{faf9f5}
\definecolor{surface-soft}{HTML}{f5f0e8}
\definecolor{surface-card}{HTML}{efe9de}
\definecolor{surface-cream-strong}{HTML}{e8e0d2}
\definecolor{surface-dark}{HTML}{181715}
\definecolor{surface-dark-elevated}{HTML}{252320}
\definecolor{surface-dark-soft}{HTML}{1f1e1b}
\definecolor{hairline}{HTML}{e6dfd8}
\definecolor{on-dark}{HTML}{faf9f5}
\definecolor{on-dark-soft}{HTML}{a09d96}
\definecolor{accent-teal}{HTML}{5db8a6}
\definecolor{accent-amber}{HTML}{e8a55a}
\definecolor{success}{HTML}{5db872}
\definecolor{warning}{HTML}{d4a017}
\definecolor{error}{HTML}{c64545}
%% alias for mdframed highlight
\definecolor{lightbg}{HTML}{f5f0e8}

\hypersetup{colorlinks=true, linkcolor=coral, urlcolor=coral}

\babelprovide[import, onchar=ids fonts]{thai}
\babelprovide[import, onchar=ids fonts]{english}

\babelfont{rm}{Noto Sans}
\babelfont[thai]{rm}{Noto Sans Thai}
\babelfont{tt}{Noto Sans Mono}
\babelfont[thai]{tt}{Noto Sans Thai}

\directlua{
  local glyph_id = node.id("glyph")
  local penalty_id = node.id("penalty")

  local function is_thai(c)
    return c >= 0x0E01 and c <= 0x0E5B
  end

  local function is_thai_combining(c)
    return (c >= 0x0E31 and c <= 0x0E3A) or (c >= 0x0E47 and c <= 0x0E4E)
  end

  luatexbase.add_to_callback("pre_linebreak_filter", function(head)
    for n in node.traverse_id(glyph_id, head) do
      local prev = n.prev
      if prev and prev.id == glyph_id and is_thai(prev.char)
         and is_thai(n.char) and not is_thai_combining(n.char) then
        local p = node.new(penalty_id)
        p.penalty = 0
        node.insert_before(head, n, p)
      end
    end
    return head
  end, "thai_break")
}

\emergencystretch=1em

\newmdenv[
  backgroundcolor=lightbg,
  linecolor=coral,
  linewidth=2pt,
  topline=false, rightline=false, bottomline=false,
  innerleftmargin=12pt, innerrightmargin=8pt,
  innertopmargin=8pt, innerbottommargin=8pt,
]{highlight}

\begin{document}

\begin{center}
  {\Large \textbf{สรุปบทความวิจัย}}\\[6pt]
  {\color{muted}\small Paper Reading Note · วิทยานิพนธ์}
\end{center}

\noindent\rule{\textwidth}{1.5pt}

%% ─── Metadata ─────────────────────────────────────────────────
\vspace{12pt}
\noindent
\begin{tabular}{@{}ll@{}}
  \textbf{ชื่อบทความ}  & Visualizing the Web Search Results with Web Search \\
                        & Visualization Using Scatter Plot \\[4pt]
  \textbf{ผู้เขียน}    & Anwar Alhenshiri, Stephen Brooks, Carolyn Watters, Michael Shepherd \\[4pt]
  \textbf{ปี / สถานที่} & 2010 · IEEE Symposium on Web Society (SWS 2010) \\[4pt]
  \textbf{arXiv / DOI}  & DOI: 10.1109/SWS.2010.5607488 \\[4pt]
  \textbf{อ่านเมื่อ}   & \today \\
\end{tabular}

\vspace{12pt}
\noindent\rule{\textwidth}{0.4pt}

%% ─── Problem ──────────────────────────────────────────────────
\section*{ปัญหาที่แก้ไข}

ระบบ Web Search ในยุคดั้งเดิมมักแสดงผลลัพธ์ในรูป \textbf{ranked list} เพียงมิติเดียว ทำให้ผู้ใช้ไม่สามารถ
สำรวจ (explore) หรือเปรียบเทียบ (compare) ผลลัพธ์ตามคุณสมบัติที่ต้องการ เช่น ความใหม่ของเนื้อหา
(freshness) หรือประเภทไฟล์ (document type) ได้อย่างมีประสิทธิภาพ:
\begin{enumerate}
  \item \textbf{ข้อจำกัดของ List-based Presentation}: การเรียงผลลัพธ์ตาม relevance rank ซ่อน
    distribution ของผลลัพธ์ตาม temporal และ format dimension ทำให้ผู้ใช้ไม่เห็นภาพรวมของ
    result space
  \item \textbf{ขาด Interactive Exploration}: ผู้ใช้ไม่มีเครื่องมือ filter หรือ navigate ผลลัพธ์ตาม
    dimension ที่ต้องการ เช่น ต้องการเฉพาะ PDF ที่เผยแพร่หลังปี 2008
\end{enumerate}

%% ─── Contribution ─────────────────────────────────────────────
\section*{ผลงานสำคัญ (Contribution)}

\begin{enumerate}[nosep, leftmargin=2em]
  \item \textbf{Dual-View Search Result Presentation}: ระบบแสดงผลลัพธ์พร้อมกันทั้ง
    traditional ranked list และ interactive 2D scatter plot view ในหน้าเดียว
  \item \textbf{2D Scatter Plot Layout}: แกน X = last modification date (temporal dimension),
    แกน Y = document type (format dimension) ทำให้ผู้ใช้เห็น distribution ของผลลัพธ์
    ได้ทันที
  \item \textbf{Linked Interactive Selection}: การ select หรือ click จุดในกราฟ scatter plot
    จะ highlight ผลลัพธ์ที่สอดคล้องใน list view และในทางกลับกัน (bidirectional linking)
\end{enumerate}

%% ─── Method ───────────────────────────────────────────────────
\section*{วิธีการ}

ระบบ Web Search Visualization ทำงานผ่าน 3 ขั้นตอนหลัก:
\begin{itemize}
  \item \textbf{Query \& Metadata Extraction}: ส่ง query ไปยัง Search Engine API,
    ดึง metadata ของแต่ละผลลัพธ์ได้แก่ last modification date และ document type
    (HTML, PDF, DOC, PPT เป็นต้น)
  \item \textbf{Dual Rendering}: นำ metadata มา render พร้อมกันใน 2 view: (1)
    traditional ranked list พร้อม title/snippet/URL และ (2) scatter plot ที่ plot
    แต่ละผลลัพธ์เป็นจุดตาม 2 แกน
  \item \textbf{Interactive Filtering}: ผู้ใช้สามารถ drag เพื่อเลือก region ใน scatter
    plot --- ผลลัพธ์ที่อยู่ใน region จะ highlight/filter ทั้งสองฝั่ง (linked views
    paradigm)
\end{itemize}

%% ─── Key Results ──────────────────────────────────────────────
\section*{ผลลัพธ์สำคัญ}

\begin{highlight}
\textbf{ข้อค้นพบหลักจาก Paper (Qualitative / Demonstration):}
\begin{itemize}[nosep]
  \item Scatter plot view ช่วยเปิดเผย \textbf{hidden patterns} ใน result set เช่น cluster
    ของ PDF วิชาการที่เผยแพร่ในช่วงปี 2005--2009 ซึ่งซ่อนอยู่ใน list view
  \item \textbf{Dual-view เสริมซึ่งกันและกัน}: ผู้ใช้ยังต้องการ ranking information จาก
    list view ดังนั้น scatter plot ทำหน้าที่ complement ไม่ใช่ replace
  \item Date และ document type เป็น metadata ที่ \textbf{extract ได้ง่าย} จาก search
    engine API และ useful สำหรับ first-pass filtering
\end{itemize}
\end{highlight}

\vspace{6pt}
\noindent\textbf{ข้อสังเกต:}
\begin{itemize}[nosep, leftmargin=2em]
  \item งานนี้เป็น \textbf{demonstration paper} ไม่มี formal user study --- ไม่มีการวัด
    task completion time, error rate หรือ user satisfaction score
  \item บริบทปี 2010 ก่อน modern search engine ที่มี rich snippets, knowledge panels
    และ AI-powered features ทำให้บางส่วนของ motivation ลดความเกี่ยวข้องลง
\end{itemize}

%% ─── Strength / Weakness ──────────────────────────────────────
\section*{จุดแข็ง / จุดอ่อน}

\noindent\textbf{จุดแข็ง:}
\begin{itemize}[nosep, leftmargin=2em]
  \item แนวคิด linked views (scatter plot $\leftrightarrow$ list) เป็น UX pattern ที่ timeless
    และยังคงถูกนำไปใช้ใน modern information visualization systems
  \item Metadata-based approach (date + type) เรียบง่ายและ generalizable --- ไม่ต้องการ
    deep semantic understanding ของเนื้อหา
  \item เป็น pioneer work ใน visual web search ที่เปิดแนวทางให้งานในภายหลัง
\end{itemize}

\vspace{6pt}
\noindent\textbf{จุดอ่อน / ข้อจำกัด:}
\begin{itemize}[nosep, leftmargin=2em]
  \item ใช้เพียง 2 dimensions (date, type) ซึ่งเป็น \textbf{shallow metadata} --- ไม่ capture
    semantic content หรือ relevance ของเนื้อหา
  \item Scatter plot coordinates ไม่แสดง \textbf{semantic proximity} --- จุดที่อยู่ใกล้กัน
    ไม่จำเป็นต้องมีเนื้อหาเกี่ยวข้องกัน
  \item ขาด empirical evidence จาก user study --- ไม่ทราบว่า dual-view ช่วยผู้ใช้จริงหรือไม่
\end{itemize}

%% ─── Relevance ────────────────────────────────────────────────
\section*{ความเกี่ยวข้องกับงานของเรา}

งานนี้มีความเกี่ยวข้องในเชิง \textbf{UX design pattern} สำหรับ information retrieval systems:
\begin{itemize}
  \item แนวคิด \textbf{dual-view (structured list + visual overview)} เป็น pattern ที่สามารถ
    นำไปประยุกต์ใช้กับระบบ browse/search scientific figures ได้ --- หาก embed รูปภาพในพื้นที่
    2D (เช่น t-SNE/UMAP บน figure embeddings) จะ meaningful กว่า metadata-based scatter plot
    มาก เพราะแสดง semantic similarity จริง
  \item สำหรับวิทยานิพนธ์: การแสดง \textbf{figure exploration interface} ที่ complement ผลการ
    ค้นหาแบบ list ด้วย visual map ของ embedding space อาจเป็น future work ที่น่าสนใจ
  \item งานนี้เตือนว่า \textbf{metadata เพียงอย่างเดียวไม่เพียงพอ} --- ระบบที่ดีควรใช้ semantic
    embeddings ในการจัด layout visual space
\end{itemize}

%% ─── Notes ────────────────────────────────────────────────────
\section*{บันทึกเพิ่มเติม}
Paper นี้เป็นตัวอย่างที่ดีของ \textit{early exploratory visualization work} ที่นำเสนอ concept
มากกว่า evidence บทเรียนสำคัญคือ: visual interface สำหรับ retrieval ต้องการ user study ที่
ชัดเจนเพื่อพิสูจน์ประโยชน์ --- อย่า assume ว่า richer visualization = better UX โดยไม่มีข้อมูล

\vspace{16pt}
\noindent\rule{\textwidth}{1pt}
\noindent{\color{muted}\small สรุปโดย Antigravity (Claude Sonnet 4.6) · \today}

\end{document}
